
\documentclass[11pt]{article}

\usepackage[a4paper,margin=1in]{geometry}
\usepackage{amsmath,amssymb,mathtools,amsthm}
\usepackage{aliascnt}
\usepackage{booktabs,array}
\usepackage{enumitem}
\usepackage{xcolor}
\usepackage{tikz}
\usetikzlibrary{arrows.meta}
\usepackage{microtype}
\usepackage[hidelinks]{hyperref}
\usepackage[nameinlink,noabbrev]{cleveref}



\hypersetup{
   pdftitle={Generating Antichains for Ideals of Prescribed Size: Universal Upper Bounds and Almost-All Lower Bounds},
  pdfauthor={L. Sunil Chandran; Rishikesh Gajjala; Kuldeep S. Meel; Daniel J. Zhang},
  pdfkeywords={ideals, downsets, antichains, model counting, inclusion-exclusion}
}

\allowdisplaybreaks
\setlist[itemize]{leftmargin=2em,itemsep=0.15em,topsep=0.25em}
\setlist[enumerate]{leftmargin=2em,itemsep=0.15em,topsep=0.25em}

\newtheorem{theorem}{Theorem}[section]

\newaliascnt{lemma}{theorem}
\newtheorem{lemma}[lemma]{Lemma}
\aliascntresetthe{lemma}

\newaliascnt{proposition}{theorem}
\newtheorem{proposition}[proposition]{Proposition}
\aliascntresetthe{proposition}

\newaliascnt{corollary}{theorem}
\newtheorem{corollary}[corollary]{Corollary}
\aliascntresetthe{corollary}

\newaliascnt{claim}{theorem}
\newtheorem{claim}[claim]{Claim}
\aliascntresetthe{claim}

\theoremstyle{definition}
\newaliascnt{definition}{theorem}
\newtheorem{definition}[definition]{Definition}
\aliascntresetthe{definition}

\newaliascnt{example}{theorem}
\newtheorem{example}[example]{Example}
\aliascntresetthe{example}

\theoremstyle{remark}
\newaliascnt{remark}{theorem}
\newtheorem{remark}[remark]{Remark}
\aliascntresetthe{remark}

\newcommand{\ID}{\operatorname{ID}}
\newcommand{\V}{\mathcal V}
\newcommand{\one}{\mathbf 1}
\newcommand{\bits}{\{0,1\}}
\newcommand{\growth}[1]{\mathfrak f(#1)}
\newcommand{\defterm}[1]{\emph{#1}}




\title {Generating Antichains for Ideals of Prescribed Size:\\Universal Upper Bounds and Almost-All Lower Bounds}

%\title{Exponentially Smaller Generating Antichains for an Ideal of Prescribed Size}
\author{L. Sunil Chandran\thanks{Department of Computer Science and Automation, Indian Institute of Science, Bengaluru, India}
\and Rishikesh Gajjala\thanks{Center for Quantum and Topological Systems, New York University Abu Dhabi, UAE}
\and Kuldeep S. Meel\thanks{Georgia Institute of Technology, Atlanta, Georgia, USA}
\and Daniel J. Zhang\footnotemark[3]}
\date{}

\begin{document}
\maketitle

\begin{abstract}
Let $\alpha(n)$ be the minimum size of a family of sets whose generated ideal
has cardinality $n$.  Chandran, Gajjala, and Meel proved that, for every
$n\ge3$,
\[
  \log(\operatorname{bl}(n)+1)\le\alpha(n)
    \le20\sqrt{\log n}\,\log\log n,
\]
where $\operatorname{bl}(n)$ is the number of maximal blocks of $1$'s in the
binary expansion of $n$.  We prove that, for every sufficiently large $n$,
\[
  \alpha(n)=O\!\left(
    \frac{(\log\log n)^2}{\log\log\log n}
  \right),
\]
substantially improving the previous universal upper bound.  We also prove
that this order is best possible for almost all positive integers: there are
absolute constants $c,C>0$ such that
\[
  c\frac{(\log\log n)^2}{\log\log\log n}
  \le\alpha(n)\le
  C\frac{(\log\log n)^2}{\log\log\log n}
\]
outside a set of natural density zero.  Consequently,
\[
  \max_{1\le m<X}\alpha(m)
  =\Theta\!\left(
    \frac{(\log\log X)^2}{\log\log\log X}
  \right).
\]

\end{abstract}

\noindent\textbf{Keywords:} Ideals, downsets, antichain, monotone decreasing families, abstract simplicial complexes, model counting, inclusion--exclusion.

\section{Introduction}

A family of sets $\mathcal F$ is called an ideal if whenever $A'\subseteq A$ and $A\in\mathcal F$, then $A'\in\mathcal F$ also.
Let $\mathcal S=\{S_1,S_2,\ldots,S_\alpha\}$ be a family of finite sets. The ideal generated by $\mathcal S$ is
\[
  \ID(\mathcal S)=2^{S_1}\cup\cdots\cup2^{S_\alpha},
\]
where $2^{S_i}$ denotes the power set of $S_i$. Let $\alpha(n)$ be the smallest number of generating sets required to obtain an ideal of size exactly $n$.

Chandran, Gajjala, and Meel~\cite{ChandranGajjalaMeel26} proved that, for every $n\ge3$,
\[
    \alpha(n)\le20\sqrt{\log n}\,\log\log n.
\]
The main upper-bound result of this paper is a substantial improvement. We first show that, if
\[
  2^{2^{q+1}}\le n<2^{2^{q+2}},
\]
then $\alpha(n)\le4q^2+3$ (Theorem~\ref{thm:first-upper}), and hence
\[
  \alpha(n)=O\bigl((\log\log n)^2\bigr).
\]
By matching whole blocks of bits, we improve this in Theorem~\ref{thm:universal-upper} to
\[
  \alpha(n)=O\!\left(\frac{(\log\log n)^2}{\log\log\log n}\right).
\]
The simpler bit-matching construction introduces the addressing idea. The proof of the sharp upper bound is given in Section~\ref{sec:universal-upper}. It uses ordinary binary addresses and a small pattern gadget to realize any prescribed residue modulo $2^B$. A direct reduction then decreases $\log\log n$ by $\log\log\log n-O(1)$ at cost $O(\log\log n)$.

For lower bounds, Chandran, Gajjala, and Meel~\cite{ChandranGajjalaMeel26} proved
\[
  \log(\operatorname{bl}(n)+1)\le\alpha(n),
\]
which already yields infinitely many values of $n$ for which $\alpha(n)\ge c\log\log n$ for some absolute constant $c>0$.
The lower-bound part of the present paper is much stronger: it matches the new universal upper bound for almost every integer.
For real $x>4$, write
\begin{equation}
  \growth{x}:=\frac{(\log\log x)^2}{\log\log\log x}.
  \label{eq:growth-function}
\end{equation}

\begin{theorem}[Main theorem]\label{thm:main}
There are absolute constants $c,C>0$ and $n_0$ such that:
\begin{enumerate}[label=\textup{(\roman*)}]
  \item $\alpha(n)\le C\growth{n}$ for every integer $n\ge n_0$;
    \item the set of integers $n\ge5$ satisfying $\alpha(n)\ge c\growth{n}$ has natural density $1$.
\end{enumerate}
Consequently,
\[
  \alpha(n)=\Theta(\growth{n})
\]
for almost all positive integers $n$, and, for every sufficiently large $X$,
\[
  \max_{1\le n<X}\alpha(n)=\Theta(\growth{X}).
\]
\end{theorem}

A pointwise lower bound of this order cannot hold for every sufficiently large integer: $\alpha(2^m)=1$ for every $m$, since the ideal $2^{[m]}$ is generated by one set. The density-one statement is therefore the appropriate typical form of a matching lower bound.

Sections~\ref{sec:bit-matching}--\ref{sec:first-range-reduction} develop the bit-matching construction and its quadratic bound. The sharp upper bound is proved in Section~\ref{sec:universal-upper}; Sections~\ref{sec:lower-counting}--\ref{sec:lower-consequences} prove the matching lower bound. The inclusion--exclusion notation introduced in Section~\ref{sec:preliminaries} is used throughout.

\subsection{Combinatorial background and motivation from computer science}

For completeness, we briefly recall the broader combinatorial context of ideals and the computer-science motivation for studying the parameter $\alpha(n)$. The material in this subsection is largely implicit in, and in part adapted from, the discussion in Chandran, Gajjala, and Meel~\cite{ChandranGajjalaMeel26}; we include it here only to make the present paper more self-contained.

Ideals (equivalently, downsets or monotone decreasing families) are among the most basic objects in extremal set theory. Together with their dual objects, namely filters or monotone increasing families, they arise naturally in several foundational results on set systems. Monotone properties, for instance, play a central role in threshold phenomena~\cite{BollobasThomason87}, and many classical inequalities and extremal arguments are most naturally formulated for ideals or dual ideals; see, for example, Kleitman's lemma~\cite{Kleitman66a} and the Four Functions Theorem of Ahlswede and Daykin~\cite{AhlswedeDaykin78}. In a number of combinatorial problems, one first reduces attention to ideals and then proves the desired extremal statement in that restricted setting. This phenomenon appears, for example, in extremal questions on intersecting families~\cite{Kleitman66b,Bollobas86} and in trace theorems, where results of Alon and Frankl show that it suffices to establish the corresponding statement for ideals~\cite{Alon83,Frankl83}. Ideals also appear under other standard names in adjacent areas: in topological combinatorics they are precisely abstract simplicial complexes~\cite{Kozlov08}, while in the language of matroids and greedoids closely related structures are studied as independence systems~\cite{KorteLovaszSchrader91}. Thus, the study of ideals lies at a natural meeting point of several themes in combinatorics.

Within this general framework, the quantity $\alpha(n)$ gives a natural extremal parameter. Equivalently, $\alpha(n)$ is the minimum size of an antichain that generates an ideal of size $n$. Although related questions concerning antichains and downsets have been studied earlier---for instance, Duffus, Howard, and Leader investigated the maximum cardinality of an antichain contained in a given ideal~\cite{DuffusHowardLeader19}---this minimization problem appears to be considerably less understood, despite its simple formulation and its structural significance.

A second, and in some sense more surprising, motivation comes from theoretical computer science, specifically from model counting. In the model-counting problem one seeks the number of satisfying assignments of a Boolean formula. Its weighted variant is a fundamental computational task with applications in probabilistic reasoning, network reliability, statistical physics, program synthesis, and verification~\cite{Roth96,SangEtAl04,DomshlakHoffmann07,XueChoiDarwiche12}. A standard approach in this area is to reduce weighted model counting to unweighted model counting, and in such reductions an important subroutine is the construction of monotone DNF formulas with exactly $n$ satisfying assignments~\cite{ChakrabortyFriedMeelVardi15}.

This task is essentially equivalent to the problem studied in this paper. Consider a monotone DNF formula on variables $x_1,\ldots,x_m$. Each term may be identified with the set of variables that occur in it, and each truth assignment may be identified with the set of variables assigned the value \textsc{false}. Under this correspondence, the satisfying assignments of the monotone DNF are in natural bijection with the members of an ideal generated by the complements of the sets corresponding to its terms; see also~\cite{ChandranGajjalaMeel26}. Consequently, the minimum number of terms in a monotone DNF having exactly $n$ satisfying assignments is precisely $\alpha(n)$.

Therefore, upper bounds on $\alpha(n)$ translate directly into upper bounds on the size of monotone DNF representations used in reductions for weighted model counting. Improvements on $\alpha(n)$ are not only of intrinsic combinatorial interest, but also have immediate algorithmic relevance through the possibility of yielding more succinct encodings in model-counting reductions.

\section{Preliminaries}\label{sec:preliminaries}

Throughout the paper, $\mathbb N=\{1,2,\ldots\}$.

\subsection{Ideals and downsets}

\begin{definition}
A collection of sets $\mathcal F$ is an ideal if, whenever $T_1\subseteq T_2$ and $T_2\in\mathcal F$, one also has $T_1\in\mathcal F$. If $\mathcal S=\{S_1,\ldots,S_s\}$ is a family of sets, then the ideal generated by $\mathcal S$ is
\[
  \ID(\mathcal S)=2^{S_1}\cup\cdots\cup2^{S_s}.
\]
The sets $S_1,\ldots,S_s$ are the generating sets of the ideal. If $\mathcal S$ is a minimal generating family, then it is an antichain.
\end{definition}

\begin{definition}
For $n\in\mathbb N$, let $\alpha(n)$ be the minimum cardinality of a family $\mathcal S$ such that $|\ID(\mathcal S)|=n$.
\end{definition}

The parameter is finite for every positive integer: $\{\varnothing\}$ generates an ideal of size $1$, and $n-1$ distinct singleton generators give an ideal of size $n$ for $n\ge2$.

For example, $\alpha(13)=3$, since
\[
  \bigl|\ID(\{\{1,2,3\},\{1,2,4\},\{3,4\}\})\bigr|=13,
\]
and two generating sets cannot produce an ideal of size $13$. Indeed, their sizes would be at most $3$; to exceed $11$, both would have size $3$, giving an ideal of size $16-2^c$ for some integer $0\le c\le3$, which is never $13$. A direct check gives $\alpha(n)\le 3$ for $1\le n\le16$, which will be enough for the base case in Section~\ref{sec:first-range-reduction}.
Throughout the paper, all logarithms are to base $2$.

\subsection{Notational conventions}

For a nonnegative integer $a$, write $\langle a\rangle$ for its binary expansion. Bits are indexed from the least significant position, so $\langle a\rangle_0$ is the least significant bit. Higher-indexed bits are taken to be $0$ if they do not occur in the finite binary expansion. For integers $u\le v$, write $[u,v]=\{u,u+1,\ldots,v\}$.

For an indexed family $\mathcal A=(A_1,\ldots,A_s)$, write
\[
  \Phi(\mathcal A):=\left|\bigcup_{i=1}^{s}2^{A_i}\right|.
\]
A nonempty set $I\subseteq[s]$ will be called an inclusion--exclusion query for $\mathcal A$, and its exponent is
\[
  e_I(\mathcal A):=\left|\bigcap_{i\in I}A_i\right|.
\]
When the family is clear, we write simply $e_I$. Since
$\bigcap_{i\in I}2^{A_i}=2^{\cap_{i\in I}A_i}$, inclusion--exclusion gives
\begin{equation}
  \Phi(\mathcal A)=
  \sum_{\varnothing\ne I\subseteq[s]}(-1)^{|I|+1}2^{e_I}.
  \label{eq:inclusion-exclusion}
\end{equation}
If the family evolves with a stage parameter $t$, we write $e_I^{(t)}$ for the exponent of query $I$ at stage $t$.

\section{Two basic lemmas}\label{sec:basic-lemmas}

\begin{lemma}[Splitting Lemma~\cite{ChandranGajjalaMeel26}]\label{lem:splitting}
For all $m,n\in\mathbb N$,
\[
  \alpha(m+n)\le\alpha(m)+\alpha(n+1).
\]
\end{lemma}

\begin{proof}
Take optimal generating families for ideals of sizes $m$ and $n+1$,
and relabel their ground sets to be disjoint. Their ideals intersect only
in the empty set. The union of the two generating families therefore
generates an ideal of size $m+(n+1)-1=m+n$.
\end{proof}

\begin{lemma}[Lifting Lemma~\cite{ChandranGajjalaMeel26}]\label{lem:lifting}
For every $t,n\in\mathbb N$,
\[
  \alpha(2^t n)\le\alpha(n).
\]
\end{lemma}

\begin{proof}
Let $\mathcal S$ generate an ideal of size $n$, and let $T$ be a set of
$t$ fresh elements. Replace each generator $S\in\mathcal S$ by $S\cup T$.
Every member of the new ideal is uniquely a union of an old ideal member
and a subset of $T$. Its cardinality is therefore $2^t n$, with no increase
in the number of generators.
\end{proof}

\section{The bit-matching construction}\label{sec:bit-matching}

We begin with a construction that matches one binary digit at a time.
It yields the explicit quadratic bound in Section~\ref{sec:first-range-reduction}
and introduces the position-addressing idea used in the block construction
of Section~\ref{sec:universal-upper}.

\begin{theorem}[Bit-matching construction]\label{thm:bit-matching}
Let $q\ge1$ be an integer and let $y$ be any integer.
Then there is an integer $m$ such that
\begin{enumerate}[label=\textup{(\roman*)}]
  \item $\alpha(m)\le2q$;
  \item $m\equiv y\pmod{2^{2^q}}$;
  \item $0<m<2^{2^{q+1}}$.
\end{enumerate}
\end{theorem}

\begin{proof}
It suffices to treat $0\le y<2^{2^q}$, since replacing $y$ by its least
nonnegative residue modulo $2^{2^q}$ does not change the desired congruence.
We construct $2q$ generating sets and correct the first $2^q$ binary digits
of the resulting ideal size, one digit at a time.

\smallskip
\noindent\textbf{The initial family.}
Let
\[
  X_0=[1,2^q+q-1].
\]
For $0\le j<q$, put
\[
  B_j=[2^j,2^{j+1}-1],
  \qquad f_j=2^q+j,
\]
and define
\begin{equation}
  S_{2j}=X_0\setminus\{f_j\},
  \qquad
  S_{2j+1}=X_0\setminus\bigl(\{f_j\}\cup B_j\bigr).
  \label{eq:bit-initial-sets}
\end{equation}
The sets $B_j$ will be useful in Claim  \ref  {clm:primary-intersection-size} to  encode bit position number $i$ by making use of the observation that 
$|B_j| = 2^j$, while the elements $f_j$ are flags
that separate the ``effect'' of a primary intersection from the effects of all its proper subsets, as you will see in Claim \ref {clm:nonprimary-larger}.


Choose fresh elements $x_0,\ldots,x_{2^q-1}$ outside $X_0$. Write
\[
  \mathcal S^{(-1)}=(S_0,\ldots,S_{2q-1})
\]
for the initial indexed family. At the end of stage $i$ we write
$\mathcal S^{(i)}=(S_0^{(i)},\ldots,S_{2q-1}^{(i)})$ and
\[
  m_i:=\Phi(\mathcal S^{(i)}).
\]
For every nonempty $I\subseteq[0,2q-1]$, we call $I$ a query at stage $t$ and write
\[
  S_I^{(t)}:=\bigcap_{r\in I}S_r^{(t)},
  \qquad e_I^{(t)}:=|S_I^{(t)}|.
\]
Thus $e_I^{(t)}$ is the exponent of the query $I$ in the inclusion--exclusion formula for $m_t$.

\smallskip
\noindent\textbf{The update at stage $i$.}
For $0\le i<2^q$, compare $\langle m_{i-1}\rangle_i$ with
$\langle y\rangle_i$. If they agree, make no change. If they disagree, add
$x_i$ to exactly one set in every pair
$\{S_{2j}^{(i-1)},S_{2j+1}^{(i-1)}\}$: add it to $S_{2j}^{(i-1)}$ when
$\langle i\rangle_j=1$, and to $S_{2j+1}^{(i-1)}$ when
$\langle i\rangle_j=0$. Define the corresponding primary query (or primary index set)
\begin{equation}
  I_i:=\{2j:\langle i\rangle_j=1\}
  \cup\{2j+1:\langle i\rangle_j=0\}.
  \label{eq:primary-index-set}
\end{equation}
Thus $|I_i|=q$, and, when an update occurs, $x_i$ is inserted precisely into
the generators indexed by $I_i$.

The next three observations explain why stage $i$ changes exactly the $i$th
unmatched bit without disturbing the lower bits.

\begin{claim}\label{clm:fresh-element-separation}
If an update occurs at stage $i$, then $x_i$ belongs to
$\bigcap_{r\in I_i}S_r^{(t)}$ for every $t\ge i$. If $j\ne i$, then
$x_i\notin\bigcap_{r\in I_j}S_r^{(t)}$ for every stage $t$.
\end{claim}

\begin{proof}
The first assertion follows directly from the update rule and from the fact
that elements are never deleted. For $j\ne i$, the binary words of $i$ and
$j$ differ in some coordinate $k$. The sets $I_i$ and $I_j$ therefore select
different members of the pair $\{2k,2k+1\}$. The element $x_i$ was inserted
into only the member selected by $I_i$, so it is absent from the member
selected by $I_j$.
\end{proof}

\begin{claim}\label{clm:primary-intersection-size}
Immediately before stage $i$,
\[
  \left|\bigcap_{r\in I_i}S_r^{(i-1)}\right|=i.
\]
\end{claim}

\begin{proof}
No fresh element $x_g$ with $g\ne i$ lies in the intersection, by Claim 
\ref {clm:fresh-element-separation}, and $x_i$ has not yet been inserted.
Every flag $f_j$ is absent from both members of its pair, so no flag survives.
For the block $B_j$, the selected generator is $S_{2j}$ precisely when
$\langle i\rangle_j=1$; in that case $B_j$ survives. When
$\langle i\rangle_j=0$, the selected generator is $S_{2j+1}$ and $B_j$ is
removed. Hence
\[
  \bigcap_{r\in I_i}S_r^{(i-1)}
  =\bigcup_{\langle i\rangle_j=1}B_j,
\]
and its cardinality is $\sum_j\langle i\rangle_j2^j=i$.
\end{proof}

\begin{claim}\label{clm:nonprimary-larger}
If $\varnothing\ne I\subsetneq I_i$, then
\[
  \left|\bigcap_{r\in I}S_r^{(i-1)}\right|>i.
\]
\end{claim}

\begin{proof}
The primary intersection is contained in the intersection indexed by $I$.
Since $I$ is a proper subset of $I_i$, for some $j$ the unique member of the
$j$th pair selected by $I_i$ is absent from $I$. Because $I\subseteq I_i$,
the other member of that pair is absent as well. The flag $f_j$, which is
missing only from those two generators, therefore survives in the
intersection indexed by $I$, but not in the primary intersection. The
containment is strict, and  Claim \ref {clm:primary-intersection-size} completes the
proof.
\end{proof}

Suppose an update occurs at stage $i$ and put
$\Delta_i=m_i-m_{i-1}$. In the inclusion--exclusion formulas for $m_i$ and
$m_{i-1}$, a query contributes to the difference exactly when it is a
nonempty subset of $I_i$. We call these the contributing queries of stage
$i$, and $I_i$ itself the primary query. Thus
\begin{equation}
  \Delta_i=
  \sum_{\varnothing\ne I\subseteq I_i}
  (-1)^{|I|+1}
  2^{e_I^{(i-1)}}.
  \label{eq:bit-increment}
\end{equation}
The primary query $I=I_i$ has exponent $i$, while every other contributing query has exponent
strictly larger than $i$. Therefore
\begin{equation}
  \Delta_i\equiv(-1)^{q+1}2^i\equiv2^i
  \pmod{2^{i+1}}.
  \label{eq:bit-increment-congruence}
\end{equation}
The last congruence holds because $-2^i\equiv2^i\pmod{2^{i+1}}$.


Thus, adding $x_i$ in the sets whose indices belong to the primary query allows us to change the bit at the $i$th position, preserving all lower bits, which have already been matched. This is referred to as the position addressing idea. It will be used in Section~\ref{sec:universal-upper}, with appropriate modification. 

We now induct on $i$. After stage $i-1$, assume that the bits below $i$ of
$m_{i-1}$ agree with those of $y$. If the $i$th bits already agree, we do
nothing. Otherwise, \eqref{eq:bit-increment-congruence} shows that the update
leaves all lower bits unchanged and toggles the $i$th bit. Consequently,
after the final stage,
\[
  m_{2^q-1}\equiv y\pmod{2^{2^q}}.
\]
Set $m=m_{2^q-1}$. The final family has $2q$ generators, so
$\alpha(m)\le2q$.

It remains to bound $m$. Initially,
\[
  |S_{2j}|=2^q+q-2,
  \qquad
  |S_{2j+1}|=2^q+q-2-2^j.
\]
For each fixed $j$, exactly half of the integers in $[0,2^q-1]$ have $j$th
bit $1$, and half have $j$th bit $0$. Hence each generator receives at most
$2^{q-1}$ fresh elements. With
\[
  A:=2^q+q-2+2^{q-1},
\]
the final sizes are at most $A$ and $A-2^j$, respectively. By the union
bound,
\[
  m\le\sum_{j=0}^{q-1}\bigl(2^A+2^{A-2^j}\bigr)
  \le\frac{3q}{2}\,2^A.
\]
For every $q\ge1$,
\[
  A+\log(3q/2)<2^{q+1}.
\]
For $q=1,2$ this follows by direct substitution; for $q\ge3$, use
$\log(3q/2)<q$ and $2q-2\le2^{q-1}$. Hence $m<2^{2^{q+1}}$. The ideal is nonempty, hence $m>0$, and the theorem
follows.
\end{proof}

\section{The first upper bound}\label{sec:first-range-reduction}

The bit-matching theorem allows us to remove a small congruent ideal size and
then divide the remaining quantity by a large power of two. Repeating this
at most twice moves an integer to the preceding double-exponential range.

\begin{theorem}\label{thm:first-upper}
Let $q\ge0$ be an integer and let $n$ be an integer satisfying
\[
  2^{2^{q+1}}\le n<2^{2^{q+2}}.
\]
Then
\[
  \alpha(n)\le4q^2+3.
\]
Consequently,
\[
  \alpha(n)=O\bigl((\log\log n)^2\bigr).
\]
\end{theorem}

\begin{proof}
We induct on $q$. For $q=0$, the range is $4\le n<16$. One or two generators realize every value in this range except $13$, while
\[
  |\ID(\{\{1,2,3\},\{1,2,4\},\{3,4\}\})|=13;
\]
hence $\alpha(n)\le3$ throughout the base range. Assume $q\ge1$.

Apply \Cref{thm:bit-matching} to $y=n+1$. We obtain $m_0$ with
\[
  0<m_0<2^{2^{q+1}},
  \qquad
  m_0\equiv n+1\pmod{2^{2^q}},
  \qquad
  \alpha(m_0)\le2q.
\]
Put
\[
  n_1:=\frac{n-m_0+1}{2^{2^q}}.
\]
By the splitting and lifting lemmas,
\begin{equation}
  \alpha(n)
  \le\alpha(m_0)+\alpha(n-m_0+1)
  \le2q+\alpha(n_1).
  \label{eq:first-reduction-one-step}
\end{equation}
If $n_1<2^{2^{q+1}}$, take $n'=n_1$; then
$\alpha(n)\le4q+\alpha(n')$.

Otherwise, apply \Cref{thm:bit-matching} once more, now to $n_1+1$, and
obtain $m_1$ with the same modulus and $\alpha(m_1)\le2q$. Put
\[
  n_2:=\frac{n_1-m_1+1}{2^{2^q}}.
\]
A second use of \Cref{lem:splitting,lem:lifting} gives
\[
  \alpha(n_1)\le2q+\alpha(n_2),
  \qquad
  \alpha(n)\le4q+\alpha(n_2).
\]
Moreover,
\[
  n_2\le\frac{n}{2^{2^{q+1}}}<2^{2^{q+1}}.
\]
Thus in either case there is a positive integer $n'<2^{2^{q+1}}$ such that
\begin{equation}
  \alpha(n)\le4q+\alpha(n').
  \label{eq:first-range-recursion}
\end{equation}
If $n'\ge4$, then $n'$ belongs to a range indexed by some $q'<q$, and the
induction hypothesis gives
\[
  \alpha(n')\le4(q-1)^2+3.
\]
The same inequality is trivial when $n'<4$. Substituting into
\eqref{eq:first-range-recursion},
\[
  \alpha(n)
  \le4q+4(q-1)^2+3
  =4q^2-4q+7
  \le4q^2+3.
\]
Since $q=\Theta(\log\log n)$, the final assertion follows.
\end{proof}

\section{Block matching and the universal upper bound}
\label{sec:universal-upper}

We now replace the one-bit updates of
Theorem~\ref{thm:bit-matching} by updates that match an entire block of
bits.  We first construct a small ideal with any prescribed residue
modulo $2^B$.  Applying this construction with $B$ slightly smaller than
$\log n$ gives a direct recurrence for $\alpha(n)$.

\subsection{A residue-matching construction}

\begin{theorem}[Block-matching construction]\label{thm:block-matching}
There are absolute constants $a,b>0$ and an integer $B_{\mathrm{res}}$ such that, for
every integer $B\ge B_{\mathrm{res}}$ and every integer $y$, there is a positive integer
$m$ satisfying
\[
  m\equiv y\pmod{2^B},
  \qquad
  \alpha(m)\le a\log B,
  \qquad
  \log m\le B+b\frac{B}{\log B}.
\]
\end{theorem}

\begin{proof}
Let $\lambda:=\log B$, and set
\begin{equation}
  r:=\left\lfloor\log\frac{B}{\log B}\right\rfloor,
  \qquad H:=2^r,
  \qquad \ell:=\left\lceil\frac{B}{H}\right\rceil,
  \qquad B^*:=H\ell.
  \label{eq:block-parameters}
\end{equation}
We take $B_{\mathrm{res}}$ sufficiently large that $r\ge1$.  These choices give
\begin{equation}
  \frac{B}{2\lambda}<H\le\frac{B}{\lambda},
  \qquad \lambda\le\ell<2\lambda+1,
  \qquad r\le \lambda,
  \qquad B\le B^*<B+H.
  \label{eq:block-parameter-bounds}
\end{equation}
Think of the $B^*$ binary positions as $H$ consecutive blocks, each of
length $\ell$.  Block $i$ occupies positions $i\ell,\ldots,(i+1)\ell-1$.
At stage $i$ we will match this block while preserving the $i\ell$ lower
bits already matched.  Thus the final ideal size will agree with $y$
modulo $2^{B^*}$, and hence modulo $2^B$.  The choice $H=2^r$ allows the
blocks to be addressed by ordinary binary words of length $r$.  Rounding
up to $B^*$ costs fewer than $H=O(B/\log B)$ extra positions, within the
size allowance in the theorem.

\smallskip
\noindent\textbf{Position and pattern generators.}
Choose pairwise disjoint sets $E_j,F_j$ for $0\le j<r$, with
\[
  |E_j|=\ell2^j,
  \qquad |F_j|=\ell,
  \qquad
  \mathsf U:=\bigcup_{j=0}^{r-1}(E_j\cup F_j).
\]
Thus
\[
  |\mathsf U|=\ell(H-1+r).
\]
The sets $E_j$ encode a block's starting position: taking their union over
the indices $j$ with $\langle i\rangle_j=1$ gives total size $i\ell$.
The flag sets $F_j$ replace the single flags in the bit-matching proof;
their size $\ell$ ensures that an inclusion--exclusion query drawn from
the generators selected at the current stage and missing a selected
position generator retains an entire flag block, placing its exponent at
or above the next block boundary.

Choose a set $T$ of size $\ell$ disjoint from $\mathsf U$, nested subsets
$T_k\subseteq T$ of size $k$ for $0\le k<\ell$, and distinct elements
$z_0,\ldots,z_{\ell-1}$ outside $\mathsf U\cup T$.  Put
\[
  \mathsf V:=T\cup\{z_0,\ldots,z_{\ell-1}\},
  \qquad |\mathsf V|=2\ell.
\]
The initial position generators are
\begin{equation}
  P_{j,1}:=(\mathsf U\setminus F_j)\cup\mathsf V,
  \qquad
  P_{j,0}:=(\mathsf U\setminus(F_j\cup E_j))\cup\mathsf V
  \quad(0\le j<r).
  \label{eq:position-generators}
\end{equation}
The initial pattern generators are
\begin{equation}
  Q_*:=\mathsf U\cup T,
  \qquad
  Q_k:=\mathsf U\cup T_k\cup\{z_k\}
  \quad(0\le k<\ell).
  \label{eq:pattern-generators}
\end{equation}
Including $\mathsf V$ in every position generator and $\mathsf U$ in
every pattern generator keeps their roles separate in a mixed
intersection.  The position generators select the part in $\mathsf U$
without removing any of the pattern in $\mathsf V$; the pattern generators
select the part in $\mathsf V$ without removing the addressed part of
$\mathsf U$.

There are $s:=2r+\ell+1$ generators.  We retain their labels as fresh
elements are added, and interpret every family below as an indexed family
of their current versions.

For $0\le i<H$, define
\[
  \mathcal P(i):=\{P_{j,\langle i\rangle_j}:0\le j<r\},
  \qquad
  C_i:=\bigcup_{\langle i\rangle_j=1}E_j.
\]
In particular, $|C_i|=i\ell$.  For
$d\in\{0,\ldots,2^\ell-1\}$, put
\[
  \mathcal Q(d):=\{Q_*\}\cup\{Q_k:\langle d\rangle_k=1\}.
\]
The next claim shows how the choice of pattern generators encodes an
arbitrary $\ell$-bit correction.

\begin{claim}[Pattern gadget]\label{clm:pattern-gadget}
For every $d\in\{0,\ldots,2^\ell-1\}$,
\begin{equation}
  Z(d):=
  \left|\ID\bigl(\{T\}\cup
       \{T_k\cup\{z_k\}:\langle d\rangle_k=1\}\bigr)\right|
  =2^\ell+d.
  \label{eq:pattern-family-Z}
\end{equation}
\end{claim}

\begin{proof}
All subsets of $T$ are present.  Each selected set
$T_k\cup\{z_k\}$ contributes exactly $2^k$ additional members containing
$z_k$, and these contributions are disjoint for distinct $k$.
\end{proof}

For example, when $\ell=3$ and $d=5=(101)_2$, select $T$ together with
$T_0\cup\{z_0\}$ and $T_2\cup\{z_2\}$.  The resulting ideal has
$8+1+4=13$ members.

\smallskip
\noindent\textbf{The update and its increment.}
Let $m_{-1}$ be the initial ideal size.  At stage $i$, suppose that
\[
  m_{i-1}\equiv y\pmod{2^{i\ell}}.
\]
This condition is vacuous at $i=0$.  Choose $d_i$ in
$\{0,\ldots,2^\ell-1\}$ so that
\begin{equation}
  (-1)^r d_i\equiv
  \frac{y-m_{i-1}}{2^{i\ell}}\pmod{2^\ell}.
  \label{eq:choose-pattern}
\end{equation}
Add one fresh element $x_i$ to precisely the generators in
$\mathcal P(i)\cup\mathcal Q(d_i)$, and denote the new ideal size by
$m_i$.  This update is performed even when $d_i=0$; the baseline
generator $Q_*$ is always selected.

As in the bit-matching proof, we first locate the primary intersection
and then identify its contribution to the increment.  Here the primary
intersection locates a block, and the selected pattern generators determine
all $\ell$ bits of the correction within that block.

\begin{claim}[Position addressing]\label{clm:position-addressing}
Immediately before stage $i$,
\begin{equation}
  \bigcap_{P\in\mathcal P(i)}P=C_i\cup\mathsf V.
  \label{eq:block-primary-intersection}
\end{equation}
\end{claim}

\begin{proof}
On the original ground set this follows from
\eqref{eq:position-generators}: every flag set is removed, and $E_j$
survives precisely when $\langle i\rangle_j=1$.  No earlier fresh element
survives.  If $g<i$, the binary addresses of $g$ and $i$ differ in some
coordinate, and $x_g$ was inserted into only one member of the
corresponding position pair.
\end{proof}

\begin{claim}[Increment from one stage]\label{clm:block-increment}
If stage $i$ inserts a fresh element into the generators in
$\mathcal P(i)\cup\mathcal Q(d)$, then the increment
$\Delta_i=m_i-m_{i-1}$ satisfies
\begin{equation}
  \Delta_i\equiv(-1)^r d\,2^{i\ell}
  \pmod{2^{(i+1)\ell}}.
  \label{eq:block-increment-congruence}
\end{equation}
\end{claim}

\begin{proof}
Every new ideal member contains $x_i$, and deleting $x_i$ gives a
bijection with the ideal generated by the old selected generators.
Consequently, inclusion--exclusion gives
\begin{equation}
  \Delta_i=
  \sum_{\varnothing\ne\mathcal J\subseteq
             \mathcal P(i)\cup\mathcal Q(d)}
    (-1)^{|\mathcal J|+1}
    2^{|\bigcap_{S\in\mathcal J}S|},
  \label{eq:block-increment-IE}
\end{equation}
where the intersections use the versions immediately before stage $i$.
We evaluate this sum modulo $2^{(i+1)\ell}$.

If $\mathcal J$ does not contain all of $\mathcal P(i)$, choose an
index $j$ whose selected position generator is missing.  Every generator
in $\mathcal J$ contains both $C_i$ and $F_j$, so the intersection has
size at least $i\ell+\ell$.  If $\mathcal J=\mathcal P(i)$, its
intersection has size $i\ell+2\ell$ by
\eqref{eq:block-primary-intersection}.  All these terms vanish modulo
$2^{(i+1)\ell}$.

The remaining queries have the form
$\mathcal P(i)\cup\mathcal L$, where
$\varnothing\ne\mathcal L\subseteq\mathcal Q(d)$.  Put
$R_*:=T$ and $R_k:=T_k\cup\{z_k\}$.  By
\eqref{eq:block-primary-intersection}, their intersections contain no
earlier fresh element and have size
\[
  i\ell+\left|\bigcap_{Q_v\in\mathcal L}R_v\right|.
\]
Substituting these queries into \eqref{eq:block-increment-IE}, and using
inclusion--exclusion once more for the pattern ideal, yields
\begin{align*}
  \Delta_i
  &\equiv
    2^{i\ell}
    \sum_{\varnothing\ne\mathcal L\subseteq\mathcal Q(d)}
      (-1)^{r+|\mathcal L|+1}
      2^{|\bigcap_{Q_v\in\mathcal L}R_v|} \\
  &=(-1)^r2^{i\ell}Z(d)
   \equiv(-1)^r d\,2^{i\ell}
     \pmod{2^{(i+1)\ell}}.
\end{align*}
This proves \eqref{eq:block-increment-congruence}.
\end{proof}

By Claim~\ref{clm:block-increment} and \eqref{eq:choose-pattern},
$m_i\equiv y\pmod{2^{(i+1)\ell}}$.  Induction through all $H$ stages
therefore gives
\[
  m:=m_{H-1}\equiv y\pmod{2^{H\ell}},
\]
and hence modulo $2^B$.  The generated ideal is nonempty, so $m>0$.

The final claim bounds both the number of generators and the extra
binary positions introduced by the construction.

\begin{claim}[Uniform size bounds]\label{clm:block-final-bounds}
For all sufficiently large $B$, the constructed ideal size $m$ satisfies
\[
  \alpha(m)\le5\log B,
  \qquad
  \log m\le B+3\frac{B}{\log B}.
\]
\end{claim}

\begin{proof}
By \eqref{eq:block-parameter-bounds},
\[
  s=2r+\ell+1\le4\lambda+2\le5\lambda
\]
for $\lambda\ge2$.  Every initial generator has size at most
\[
  |\mathsf U|+\ell=\ell(H+r)=B^*+r\ell.
\]
Each generator receives at most $H$ fresh elements.  Thus the largest
final generator size $M$ satisfies
\[
  M\le B^*+r\ell+H
   \le B+\frac{2B}{\lambda}+2\lambda^2+\lambda.
\]
The union bound now gives
\[
  \log m\le M+\log s
  \le B+\frac{2B}{\lambda}+2\lambda^2+\lambda+\log(4\lambda+2).
\]
Since $B=2^\lambda$, increasing the absolute threshold $B_{\mathrm{res}}$ ensures
\[
  2\lambda^2+\lambda+\log(4\lambda+2)\le\frac{B}{\lambda}
\]
for all $B\ge B_{\mathrm{res}}$.  Hence $\alpha(m)\le5\log B$ and
$\log m\le B+3B/\log B$.
\end{proof}

The congruence, positivity, and Claim~\ref{clm:block-final-bounds} prove
the theorem with fixed absolute constants $a=5$ and $b=3$.
\end{proof}

\subsection{Direct reduction}

To use the splitting lemma, we match the residue $n+1$: this makes
$n-m+1$ divisible by the chosen power of two.  Choosing $B$ close to
$\log n$ bounds the quotient's logarithm by
$O(\log n/\log\log n)$, at a cost of $O(\log\log n)$ generators.

\begin{lemma}[Direct reduction]\label{lem:direct-reduction}
There are absolute constants $a,C_1>0$ and an integer $n_{\mathrm{red}}$ such that,
for every integer $n\ge n_{\mathrm{red}}$, there is an integer $1\le n'<n$ with
\[
  \alpha(n)\le a\log\log n+\alpha(n'),
  \qquad
  \log n'\le C_1\frac{\log n}{\log\log n}.
\]
\end{lemma}

\begin{proof}
Fix $a,b,B_{\mathrm{res}}$ from Theorem~\ref{thm:block-matching}, and then fix
$A:=4b+4$.  For $L:=\log n$ and $\tau:=\log L=\log\log n$, set
\[
  B:=\left\lfloor L\left(1-\frac{A}{\tau}\right)\right\rfloor.
\]
Choose $n_{\mathrm{red}}$ large enough that $\tau\ge\max\{2A,4\}$,
$B\ge B_{\mathrm{res}}$, and $L\ge4$.  Then
$B\ge L/2-1\ge L/4$, so $\log B\ge \tau-2$.  Apply
Theorem~\ref{thm:block-matching} with $y=n+1$.  Its representative $m$
satisfies
\[
  \alpha(m)\le a\log B\le a\tau,
\]
and
\[
  \log m
  \le B+b\frac{B}{\log B}
  \le L-\frac{AL}{\tau}+\frac{bL}{\tau-2}
  \le L-\frac{(A-2b)L}{\tau}<L.
\]
Thus $m<n$, and
\[
  n':=\frac{n-m+1}{2^B}
\]
is a positive integer.  Since $m\ge1$ and $B\ge1$, we have
$n'\le n/2^B<n$.  The splitting and lifting lemmas give
\[
  \alpha(n)
  \le\alpha(m)+\alpha(n-m+1)
  \le a\tau+\alpha(n').
\]
Finally,
\[
  \log n'\le L-B\le\frac{AL}{\tau}+1
  \le\frac{(A+1)L}{\tau},
\]
where $L/\tau\ge1$ for the chosen range.  Take $C_1=A+1$.
\end{proof}

\subsection{The universal upper bound}

Writing $\tau=\log\log n$, each reduction costs $O(\tau)$ and decreases $\tau$ by
at least $\log \tau-O(1)$ until a fixed finite range is reached.  The function
$P(\tau)=\tau^2/\log \tau$ measures this total cost: over an interval of length
comparable to $\log \tau$, its increase is of order $\tau$.  The induction below
makes this comparison precise.

\begin{theorem}[Universal upper bound]\label{thm:universal-upper}
There is an absolute constant $C>0$ such that, for every integer $n\ge5$,
\[
  \alpha(n)\le
  C\frac{(\log\log n)^2}{\log\log\log n}.
\]
\end{theorem}

\begin{proof}
For real $u\ge4$, put $P(u):=u^2/\log u$.  Its derivative satisfies
\[
  P'(u)=\frac{2u}{\log u}
         -\frac{u}{(\log u)^2\ln2}
       \ge\frac{u}{\log u}>0.
\]
For $\tau\ge8$, every $u\in[\tau-\tfrac12\log \tau,\tau]$ is at least $\tau/2$ and
at least $4$.  Integrating the preceding inequality gives
\begin{equation}
  P(\tau)-P\!\left(\tau-\tfrac12\log \tau\right)
  \ge\frac{\log \tau}{2}\,\frac{\tau}{2\log \tau}
  =\frac \tau4.
  \label{eq:upper-potential-gap}
\end{equation}

Let $a,C_1,n_{\mathrm{red}}$ be as in Lemma~\ref{lem:direct-reduction}.
Choose an integer $n_0\ge n_{\mathrm{red}}$ so large that every $n\ge n_0$, with
$\tau:=\log\log n$, satisfies $\tau\ge8$ and
$\log \tau\ge2\max\{\log C_1,0\}$.  Set
\[
  F_0:=\max_{1\le k<n_0}\alpha(k),
  \qquad D:=4a.
\]
The finite maximum $F_0$ exists because the elementary construction in
the preliminaries gives $\alpha(k)\le\max\{1,k-1\}$ for every positive
integer $k$.  We prove by strong induction on integers $n\ge n_0$ that
\begin{equation}
  \alpha(n)\le F_0+D\,P(\log\log n).
  \label{eq:upper-potential-induction}
\end{equation}
Apply Lemma~\ref{lem:direct-reduction}, obtaining $1\le n'<n$.
If $n'<n_0$, then
\[
  \alpha(n)\le a\tau+F_0\le F_0+D\,P(\tau),
\]
because $P(\tau)\ge \tau$ for $\tau\ge8$.

If $n'\ge n_0$, define $\tau':=\log\log n'$.  The reduction lemma gives
\[
  \tau'\le \tau-\log \tau+\log C_1
      \le \tau-\tfrac12\log \tau.
\]
Here $\tau'\ge\log\log n_0\ge8$, so $P(\tau')$ is defined and lies in the
range where $P$ is increasing.  By the induction hypothesis and
\eqref{eq:upper-potential-gap},
\[
  \alpha(n)
  \le a\tau+F_0+D\,P(\tau')
  \le F_0+D\,P(\tau),
\]
since $D/4=a$.  This proves \eqref{eq:upper-potential-induction}.

Thus $\alpha(n)=O(\growth{n})$ as $n\to\infty$.
Since $\growth{n}>0$ for every integer $n\ge5$, increasing the absolute
constant covers the finitely many values $5\le n<n_0$ and proves the
stated bound.
\end{proof}

\section{The lower bound as a counting problem}\label{sec:lower-counting}

Theorem~\ref{thm:universal-upper} proves the universal upper bound in
Theorem~\ref{thm:main}. To obtain a matching lower bound for almost all
integers, we count the ideal sizes that can be generated by a small family.
For integers $B,K\ge1$, let
\[
  \V(B,K):=\{\Phi(\mathcal A):
       \mathcal A=(A_1,\ldots,A_s),\quad
       1\le s\le K,\quad |A_i|\le B\text{ for all }i\}.
\]
\begin{theorem}[Counting attainable ideal sizes]\label{thm:attainable-count}
There exists an absolute integer $B_{\mathrm{cnt}}$ such that, for all integers
$B\ge B_{\mathrm{cnt}}$ and $K$ satisfying
\[
  1\le K\le\frac{(\log B)^2}{100\log\log B},
\]
one has
\[
  |\V(B,K)|\le2^{B-\sqrt B}.
\]
\end{theorem}

For the proof, take $B\ge2^{20}$ and introduce the auxiliary parameters
\begin{equation}
  \lambda:=\log B,\qquad
  h:=\left\lfloor\frac{\lambda}{4}\right\rfloor,\qquad
  J:=\left\lfloor\frac Kh\right\rfloor,\qquad
  \Gamma:=\frac{B}{\lambda^{6J}}.
  \label{eq:main-parameters}
\end{equation}
We will prove the stronger estimate $|\V(B,K)|\le2^{B-\Gamma/4}$ and
then check that $\Gamma/4\ge\sqrt B$ for all sufficiently large $B$,
uniformly over the stated range of $K$.

We may pad every generating family to a $K$-tuple by repeating its final
generator. Delete ground elements that lie in no generator, and identify the
remaining ground set with $[N]$. Then
\begin{equation}
  N=\left|\bigcup_{i=1}^K A_i\right|
    \le\sum_{i=1}^K|A_i|\le KB.
  \label{eq:ground-bound}
\end{equation}
For every nonempty $I\subseteq[K]$, write
$e_I=|\bigcap_{i\in I}A_i|$. Inclusion--exclusion gives
\begin{equation}
  \Phi(\mathcal A)
  =\sum_{\varnothing\ne I\subseteq[K]}(-1)^{|I|+1}2^{e_I}.
  \label{eq:IE}
\end{equation}

Recording all $2^K-1$ exponents independently gives at most
$(B+1)^{2^K-1}$ possible vectors. Near the upper end of the permitted range
of $K$, this is much larger than $2^B$, so it cannot prove the desired
saving. We instead use the overlap structure to determine many large terms
from a short common record, and count the remaining sum by its numerical
range.

Section~\ref{sec:headers} selects two overlap scales and partitions the
generator indices into clusters and small residual components. Only the
scale index and this partition form the structural header. For a fixed
header, Section~\ref{sec:fixed-header} finds a small exceptional ground set
$U$ outside which each cluster agrees with one representative. The incidence
record on $U$ and a short transcript of intersection exponents on its
complement determine the entire central sum $Z$. Every noncentral term has
exponent at most $B-\rho$, so the sum of these terms lies in one short
integer interval. We multiply these three counts within each header and
then sum over headers; Figure~\ref{fig:counting-scheme} summarizes the
calculation.

\begin{figure}[htbp]
\centering
\begin{tikzpicture}[
  box/.style={draw, rounded corners=2pt, align=center,
              inner xsep=4pt, inner ysep=6pt, font=\small},
  flow/.style={-{Latex[length=2mm]}, thick}
]
  \node[box, text width=7cm] (header) at (0,0) {
    \textbf{Structural header: scale index and partition}\\[0.2em]
    at most $(J+1)K^K$ choices
  };
  \node[box, text width=3.9cm, minimum height=1.5cm] (record) at (-4.8,-2.3) {
    Incidence record on $U$\\[0.4em]
    $(KB+1)2^{2K^2\rho/\lambda^6}$
  };
  \node[box, text width=3.9cm, minimum height=1.5cm] (transcript) at (0,-2.3) {
    Central exponent transcript\\[0.4em]
    $(B+1)^{2^J(1+K2^h)}$
  };
  \node[box, text width=3.9cm, minimum height=1.5cm] (interval) at (4.8,-2.3) {
    Integer interval for $\Phi(\mathcal A)-Z$\\[0.4em]
    $2^{B-\rho+K+2}$
  };
  \node at (-2.4,-2.3) {$\times$};
  \node at (2.4,-2.3) {$\times$};
  \draw[thick] (header.south) -- (0,-1.2);
  \draw[thick] (-4.8,-1.2) -- (4.8,-1.2);
  \draw[flow] (-4.8,-1.2) -- (record.north);
  \draw[flow] (0,-1.2) -- (transcript.north);
  \draw[flow] (4.8,-1.2) -- (interval.north);
\end{tikzpicture}
\caption{For a fixed header with selected coarse scale $\rho$, the three
lower factors bound the number of attainable values. Summing over headers
completes the count. The sets and records are defined below.}
\label{fig:counting-scheme}
\end{figure}

\section{Structural headers}\label{sec:headers}

For $\theta>0$, define the proximity graph $G_\theta$ on $[K]$ by
\[
  ij\in E(G_\theta)
  \quad\Longleftrightarrow\quad
  |A_i\cap A_j|>B-\theta.
\]
Thus $G_{\theta'}\subseteq G_\theta$ whenever $\theta'\le\theta$. Consider
the scales
\[
  \rho_j:=\frac{B}{\lambda^{6j}}
  \qquad(0\le j\le J+1).
\]
Let $\nu_j$ be the maximum number of pairwise disjoint connected
$h$-vertex sets in $G_{\rho_j}$. These packing numbers satisfy
\[
  J\ge\nu_0\ge\nu_1\ge\cdots\ge\nu_{J+1}\ge0.
\]
Among the $J+1$ adjacent pairs, at least one has equal entries. Choose such
an index $j\in\{0,\ldots,J\}$, and put
\begin{equation}
  \rho:=\rho_j,\qquad
  \rho':=\rho_{j+1}=\rho/\lambda^6.
  \label{eq:selected-scales}
\end{equation}
In particular, $\rho\ge\rho_J=\Gamma$.

Write $r:=\nu_j=\nu_{j+1}$ and choose a maximum packing
$C_1,\ldots,C_r$ of connected $h$-vertex sets in the fine graph
$G_{\rho'}$. These sets also form a maximum packing in the coarse graph
$G_\rho$. Let
\[
  R:=[K]\setminus\bigcup_{a=1}^r C_a,
\]
and let $H_1,\ldots,H_p$ be the connected components of $G_\rho[R]$.
Every $H_b$ has fewer than $h$ vertices. Indeed, a connected graph with at
least $h$ vertices contains a connected $h$-vertex set: take a spanning tree
and repeatedly delete leaves. Such a set in $R$ would extend the coarse
packing, contradicting its maximality.

The two scales serve different purposes. Fine connectivity will give an
exceptional set of size at most $2K\rho'$, whose incidence words cost at
most $2K^2\rho'$ bits to record. A coarse nonedge gives a loss of $\rho$
in the full intersection exponent. Using a single scale would give a
recording cost as large as $2K^2\rho$ against a saving of only $\rho$,
so these estimates would not yield a net saving. Instead, the choice
$\rho'=\rho/\lambda^6$, together with $K\le \lambda^2$, gives
\[
  2K^2\rho'\le\frac{2\rho}{\lambda^2}=o(\rho).
\]
The exponent $6$ is a convenient choice with room to spare, rather than a
sharp threshold. We start at $\rho_0=B$. What matters is the separation
between adjacent scales and that the smallest coarse scale remains large
enough to absorb the other recording costs, as checked in
Section~\ref{sec:lower-consequences}.

We consider several thresholds so that both requirements can be met by
the same partition. A maximum packing at a smaller threshold need not
be maximum at the preceding larger threshold. After removing its clusters,
the graph at the larger threshold could still contain a component of size
at least $h$. However, the integer packing number can decrease at most $J$
times, so some consecutive pair of thresholds has equal packing numbers.
A maximum packing at the smaller threshold is then also maximum at the
larger threshold. Its clusters are connected at the smaller threshold,
while the remaining components are small even at the larger threshold.
Only the packing numbers need to agree; the graphs themselves may differ.

We call the pair consisting of $j$ and the set partition
\[
  \{C_1,\ldots,C_r,H_1,\ldots,H_p\}
\]
a \defterm{structural header}. Blocks of size $h$ are precisely the
clusters $C_a$, and all smaller blocks are residual components $H_b$; hence
no separate block labels are needed. Order the clusters and the residual
components by their least elements, and choose the canonical root
$d_a:=\min C_a$ in each cluster. A family is consistent with a header if the
header can be obtained from it by the preceding construction.

Every family admits at least one header. There are at most $J+1$ choices
for $j$ and at most $K^K$ set partitions of $[K]$: a partition is encoded by
the function sending each vertex to the least element of its block. Thus
the number of headers is at most
\begin{equation}
  (J+1)K^K.
  \label{eq:header-count}
\end{equation}
The graphs and spanning trees need not be recorded. Their role is to
certify the properties of the partition used below.

\section{Values within one structural header}\label{sec:fixed-header}

Fix a header and consider a family consistent with it. Define the
exceptional set intrinsically by
\[
  U:=\{x\in[N]:\text{membership of }x\text{ is nonconstant on some }C_a\},
  \qquad u:=|U|,\qquad W:=[N]\setminus U.
\]
Equivalently,
$U=\bigcup_a\bigcup_{i\in C_a}(A_i\triangle A_{d_a})$.
For any spanning tree of a cluster, an element has nonconstant membership
on that cluster if and only if its membership differs across some tree
edge. Choose, for the purpose of bounding $u$, a spanning tree in each fine
graph $G_{\rho'}[C_a]$. Every edge $ij$ of such a tree satisfies
\[
  |A_i\triangle A_j|
  =|A_i|+|A_j|-2|A_i\cap A_j|<2\rho'.
\]
Taking the union over the $r(h-1)$ tree edges gives
\begin{equation}
  u\le2r(h-1)\rho'\le\frac{2K\rho}{\lambda^6}.
  \label{eq:u-bound}
\end{equation}
By the definition of $W$, all generators in a cluster agree there:
\begin{equation}
  A_i\cap W=A_{d_a}\cap W\qquad(i\in C_a).
  \label{eq:cluster-agreement}
\end{equation}

\subsection{Recording the central terms}

For a nonempty inclusion--exclusion query $I\subseteq[K]$, put
\[
  S(I):=\{a\in[r]:I\cap C_a\ne\varnothing\},
  \qquad Q(I):=I\cap R.
\]
Call $I$ \defterm{central} if $Q(I)$ is empty or is contained in one
residual component $H_b$. This class is deliberately generous: some central
queries may already have small exponents, but recording them all avoids a
finer classification. There is no restriction on how many clusters $I$
meets. On $W$, only their subset $S(I)$ matters, and there are at most
$2^r$ choices for it. For every pair
\[
  S\subseteq[r],\quad Q\subseteq R,\quad
  (S,Q)\ne(\varnothing,\varnothing),
\]
with $Q$ empty or contained in one $H_b$, record the exponent
\[
  w_{S,Q}:=
  \left|W\cap\bigcap_{a\in S}A_{d_a}\cap\bigcap_{i\in Q}A_i\right|.
\]
An intersection over an empty index set is understood to be $[N]$.
Every recorded exponent lies in $\{0,\ldots,B\}$. Since $|H_b|<h$ and
$r\le J$, the number $M_{\mathrm{cen}}$ of recorded exponents satisfies
\begin{equation}
  M_{\mathrm{cen}}\le2^r\left(1+\sum_{b=1}^p(2^{|H_b|}-1)\right)
    \le2^J(1+K2^h).
  \label{eq:center-query-count}
\end{equation}
In a canonical order determined by the header, these exponents form a
transcript with at most $(B+1)^{2^J(1+K2^h)}$ possible values.

We also record $u$ and the incidence words
\[
  \sigma(x):=(\one_{x\in A_1},\ldots,\one_{x\in A_K})\in\bits^K
  \qquad(x\in U),
\]
listed in lexicographic order, with repetitions. For each query $I$, this
record determines
\[
  u_I:=\left|U\cap\bigcap_{i\in I}A_i\right|.
\]
For a central query, agreement on $W$ gives
\[
  e_I=u_I+w_{S(I),Q(I)}.
\]
Consequently, the incidence record and the transcript together determine
the exact integer
\[
  Z:=\sum_{\substack{\varnothing\ne I\subseteq[K]\\ I\text{ central}}}
          (-1)^{|I|+1}2^{e_I}.
\]
Thus even if many original queries are central, their entire signed sum
is fixed by these two records. This is an exact integer reconstruction;
there is no need to track binary carries or to record each central term.

There are at most $KB+1$ choices for $u$. For each choice, the number of
lists of $u$ incidence words is at most $2^{Ku}$. By
\eqref{eq:u-bound}, the total number of possible incidence records is
therefore at most
\begin{equation}
  (KB+1)2^{2K^2\rho/\lambda^6}.
  \label{eq:exceptional-record-count}
\end{equation}

\subsection{Bounding all other terms together}

If $I$ is noncentral, its residual part contains indices $i,j$ from two
different components of $G_\rho[R]$. These indices are not adjacent in
$G_\rho$, so
\[
  e_I\le|A_i\cap A_j|\le B-\rho.
\]
We do not record these remaining exponents. Their common loss of $\rho$
allows us to bound the sum of all their terms at once. There are fewer than
$2^K$ original queries, so the aggregate noncentral contribution satisfies
\begin{equation}
  \left|\Phi(\mathcal A)-Z\right|
  \le2^{B-\rho+K}.
  \label{eq:tail-final-bound}
\end{equation}
For each fixed incidence record and transcript, $Z$ is fixed, and
$\Phi(\mathcal A)-Z$ is an integer. Since $\rho\le B$, the interval in
\eqref{eq:tail-final-bound} contains at most $2^{B-\rho+K+2}$ integers.
Combining the record, transcript, and interval counts, the number of values
arising within this header is at most
\begin{equation}
  (KB+1)2^{2K^2\rho/\lambda^6}
  (B+1)^{2^J(1+K2^h)}
  2^{B-\rho+K+2}.
  \label{eq:fixed-header-count}
\end{equation}

\section{The counting theorem and its consequences}\label{sec:lower-consequences}

\subsection{Proof of the counting theorem}

We now check that the information recorded above costs less than the saving
from the noncentral terms. All estimates in this subsection are uniform
over $1\le K\le \lambda^2/(100\log \lambda)$.
For $\lambda\ge20$,
\[
  h\ge \lambda/5,\qquad
  J\le K/h\le\frac{\lambda}{20\log \lambda}.
\]
Consequently,
\begin{equation}
  \Gamma=\frac{2^\lambda}{\lambda^{6J}}
     \ge2^{7\lambda/10}.
  \label{eq:Gamma-lower}
\end{equation}
Put
\[
  d:=\frac{2K^2}{\lambda^6}
\]
and
\begin{align*}
  L_{\mathrm{rec}}:={}&\log(J+1)+K\log K+\log(KB+1)\\
             &+2^J(1+K2^h)\log(B+1)+K+2.
\end{align*}
Using $K\le \lambda^2$, we have $d\le2/\lambda^2=o(1)$. The transcript term satisfies
\[
  2^J(1+K2^h)\log(B+1)
  \le2^{\lambda/4+\lambda/(20\log \lambda)+1}\lambda^2(\lambda+1).
\]
All other terms in $L_{\mathrm{rec}}$ are $O(\lambda^2\log \lambda)$. Dividing these bounds by
the lower bound in \eqref{eq:Gamma-lower} shows that
\[
  \frac{L_{\mathrm{rec}}}{\Gamma}\longrightarrow0
  \qquad(\lambda\longrightarrow\infty)
\]
uniformly in the permitted range of $K$: the exponential factor in the
transcript bound has exponent $\lambda/4+o(\lambda)$, whereas that in
\eqref{eq:Gamma-lower} has exponent $7\lambda/10$.
Thus there is an absolute integer $\lambda_0\ge20$ such that, for all $\lambda\ge \lambda_0$
and all admissible $K$,
\[
  d\le\frac14,\qquad L_{\mathrm{rec}}\le\frac\Gamma2.
\]
Moreover, \eqref{eq:Gamma-lower} gives
\[
  \frac{\Gamma}{4\sqrt B}
  \ge\frac{2^{\lambda/5}}4\longrightarrow\infty.
\]
Increasing $\lambda_0$ if necessary therefore ensures $\Gamma/4\ge\sqrt B$
throughout the same range.

Every eligible family has a header, and counting the same value more than
once only increases an upper bound. Sum \eqref{eq:fixed-header-count} over
the at most $(J+1)K^K$ headers in \eqref{eq:header-count}. Since every
selected scale satisfies $\rho\ge\Gamma$ and $1-d>0$, this gives
\[
  \log|\V(B,K)|
  \le B-(1-d)\Gamma+L_{\mathrm{rec}}
  \le B-\frac{3\Gamma}{4}+\frac\Gamma2
  =B-\frac\Gamma4
  \le B-\sqrt B.
\]
Taking $B_{\mathrm{cnt}}:=2^{\lambda_0}$ proves Theorem~\ref{thm:attainable-count}.



\subsection{Typical and almost-all lower bounds}

The exact count has a direct probabilistic interpretation.

\begin{corollary}[Typical lower bound in one dyadic block]\label{cor:block-probability}
Under the assumptions of Theorem~\ref{thm:attainable-count}, let
\[
  H_B=[2^{B-1},2^B)\cap\mathbb Z
\]
and choose $Y_B$ uniformly from $H_B$. Then
\[
  \Pr\bigl(\alpha(Y_B)>K\bigr)\ge 1-2^{1-\sqrt B}.
\]
\end{corollary}

\begin{proof}
If $n\in H_B$ and $\alpha(n)\le K$, take a witnessing generating family
$A_1,\ldots,A_s$, where $s\le K$. Since each $2^{A_i}$ is contained in
the ideal,
\[
  2^{|A_i|}\le n<2^B,
\]
so $|A_i|\le B-1$. Hence $n\in\V(B,K)$. Theorem~\ref{thm:attainable-count}
therefore shows that at most $2^{B-\sqrt B}$ of the $2^{B-1}$ integers in
$H_B$ have $\alpha(n)\le K$. Dividing by $2^{B-1}$ proves the claim.
\end{proof}

\begin{corollary}[Almost-all order]\label{cor:almost-all-order}
There are absolute constants $c,C>0$ such that
\[
  \lim_{X\to\infty}\frac1X
    \left|\left\{5\le n\le X:
  c\frac{(\log\log n)^2}{\log\log\log n}
  \le \alpha(n)\le
  C\frac{(\log\log n)^2}{\log\log\log n}
  \right\}\right|=1.
\]
\end{corollary}

\begin{proof}
For every sufficiently large integer $B$, put
\[
  \lambda_B:=\log B,
  \qquad
  K_B:=\left\lfloor\frac{\lambda_B^2}{200\log \lambda_B}\right\rfloor,
\]
and define
\[
  \mathcal E_B:=\{n\in H_B:\alpha(n)\le K_B\}.
\]
By \Cref{cor:block-probability},
\begin{equation}
  \frac{|\mathcal E_B|}{|H_B|}
  \le\varepsilon_B:=2^{1-\sqrt B}.
  \label{eq:block-exceptional-proportion}
\end{equation}
In particular, $\varepsilon_B\to0$ as $B\to\infty$.

For $n\in H_B$,
\[
  B-1\le\log n<B.
\]
Hence, with $\tau:=\log\log n$,
\begin{equation}
  \lambda_B-1\le \tau<\lambda_B.
  \label{eq:block-loglog-comparison}
\end{equation}
For all sufficiently large $B$, this gives
\[
  \growth{n}=\frac{\tau^2}{\log \tau}
  \le\frac{2\lambda_B^2}{\log \lambda_B}.
\]
Also, after increasing the threshold if necessary,
\begin{equation}
  K_B\ge\frac{\lambda_B^2}{400\log \lambda_B}
  \ge\frac1{800}\growth{n}
  \qquad(n\in H_B).
  \label{eq:KB-versus-growth}
\end{equation}
Thus every $n\in H_B\setminus \mathcal E_B$ satisfies
\[
  \alpha(n)>K_B\ge\frac1{800}\growth{n}.
\]

Let
\[
  \mathcal E:=\left\{n:\alpha(n)<\frac1{800}\growth{n}\right\},
\]
ignoring the finitely many integers for which \eqref{eq:growth-function} has
not yet been defined. By \eqref{eq:KB-versus-growth},
$\mathcal E\cap H_B\subseteq \mathcal E_B$ for every sufficiently large $B$.

We show that $\mathcal E$ has natural density $0$. Given $\delta>0$, choose $B_1$ so
large that all preceding estimates, in particular
\eqref{eq:KB-versus-growth}, hold and $\varepsilon_B\le\delta$ for every
$B\ge B_1$. Let $C_{B_1}$ be
the number of exceptional integers lying in the finitely many blocks with
$B<B_1$. If $2^{M-1}\le X<2^M$, then
\begin{align*}
  |\mathcal E\cap[1,X]|
  &\le C_{B_1}+\sum_{B=B_1}^{M}|\mathcal E_B|\\
  &\le C_{B_1}+\delta\sum_{B=B_1}^{M}2^{B-1}\\
  &\le C_{B_1}+\delta 2^M
  \le C_{B_1}+2\delta X.
\end{align*}
Dividing by $X$, taking the limit superior, and then letting $\delta\downarrow0$
shows that $\mathcal E$ has density $0$. Therefore
\[
  \alpha(n)\ge\frac1{800}\growth{n}
\]
for almost all $n$. The universal upper bound in Theorem~\ref{thm:universal-upper}
holds for every sufficiently large $n$. Combining the two statements proves the corollary.
\end{proof}

\subsection{The extremal consequence}

Put
\[
  M(X):=\max_{1\le n<X}\alpha(n).
\]
The blockwise form of the lower bound gives the extremal conclusion directly.
Let $X$ be sufficiently large and put
\[
  B:=\lfloor\log X\rfloor,
  \qquad \lambda:=\log B,
  \qquad
  K:=\left\lfloor\frac{\lambda^2}{200\log \lambda}\right\rfloor.
\]
Then $2^B\le X$, and the whole block
\[
  H_B=[2^{B-1},2^B)\cap\mathbb Z
\]
lies below $X$.  For large $X$, the error term in
\Cref{cor:block-probability} is smaller than $1$, so some $n\in H_B$
satisfies $\alpha(n)>K$.  Hence
\[
  M(X)\ge K.
\]
Since $B=\Theta(\log X)$, we have
\[
  \lambda=\log B=\Theta(\log\log X),
  \qquad
  \log \lambda=\Theta(\log\log\log X),
\]
and therefore
\[
  M(X)=\Omega\!\left(
    \frac{(\log\log X)^2}{\log\log\log X}
  \right).
\]

For the reverse inequality, \Cref{thm:universal-upper} gives
\[
  \alpha(n)\le C\frac{(\log\log n)^2}{\log\log\log n}
\]
for every sufficiently large $n$.  The comparison function on the right is
increasing for all sufficiently large arguments.  Enlarging $C$ to absorb
the finitely many smaller values gives
\[
  M(X)=O\!\left(
    \frac{(\log\log X)^2}{\log\log\log X}
  \right).
\]
Together with the lower bound, this proves the extremal assertion in
\Cref{thm:main}.

\begin{remark}
The counting theorem uses only the upper bounds $|A_i|\le B$. It requires no
balancedness assumption, no restriction on intersections of intermediate
orders, and no separate bound on the ambient ground set; the latter follows
from \eqref{eq:ground-bound}.
\end{remark}

\section*{AI assistance and formalization}

The authors initially obtained the universal upper bound
$O((\log\log n)^2)$ without any AI assistance. Subsequent discussions with
OpenAI's GPT-5.5 contributed to improving this bound to
$O((\log\log n)^2/\log\log\log n)$. The matching almost-all
lower bound was developed through interaction
between the authors and OpenAI's GPT-6.

The universal upper bound and a matching lower bound at arbitrarily large
inputs have also been formalized in Lean, with assistance from multiple
AI models. The formalization is available at
\url{https://github.com/rishigajjala/antichain-of-given-size}.
The authors take full responsibility
for the mathematical statements, proofs, formalizations,
and final presentation.

\begin{thebibliography}{99}
\bibitem{AhlswedeDaykin78}
R.~Ahlswede and D.~E. Daykin,
\emph{An inequality for the weights of two families of sets, their unions and intersections},
Z. Wahrscheinlichkeitstheorie verw. Gebiete \textbf{43} (1978), 183--185.

\bibitem{Alon83}
N.~Alon,
\emph{On the density of sets of vectors},
Discrete Math. \textbf{46} (1983), 199--202.

\bibitem{Bollobas86}
B.~Bollob\'as,
\emph{Combinatorics: Set Systems, Hypergraphs, Families of Vectors, and Combinatorial Probability},
Cambridge University Press, 1986.

\bibitem{BollobasThomason87}
B.~Bollob\'as and A.~G. Thomason,
\emph{Threshold functions},
Combinatorica \textbf{7} (1987), 35--38.

\bibitem{ChakrabortyFriedMeelVardi15}
S.~Chakraborty, D.~Fried, K.~S. Meel, and M.~Y. Vardi,
\emph{From weighted to unweighted model counting},
in \emph{Proceedings of the Twenty-Fourth International Joint Conference on Artificial Intelligence},
AAAI Press, 2015, pp.~689--695.

\bibitem{ChandranGajjalaMeel26}
L.~Sunil Chandran, R.~Gajjala, and K.~S. Meel,
\emph{On the smallest antichain that generates an ideal of a given size},
J. Combin. Theory Ser. A \textbf{222} (2026), 106182,
doi:10.1016/j.jcta.2026.106182.

\bibitem{DomshlakHoffmann07}
C.~Domshlak and J.~Hoffmann,
\emph{Probabilistic planning via heuristic forward search and weighted model counting},
J. Artificial Intelligence Res. \textbf{30} (2007), 565--620.

\bibitem{DuffusHowardLeader19}
D.~Duffus, D.~Howard, and I.~Leader,
\emph{The width of downsets},
European J. Combin. \textbf{79} (2019), 46--59.

\bibitem{Frankl83}
P.~Frankl,
\emph{On the trace of finite sets},
J. Combin. Theory Ser. A \textbf{34} (1983), 41--45.

\bibitem{Kleitman66a}
D.~J. Kleitman,
\emph{Families of non-disjoint subsets},
J. Combin. Theory \textbf{1} (1966), 153--155.

\bibitem{Kleitman66b}
D.~J. Kleitman,
\emph{On a combinatorial conjecture of Erd\H{o}s},
J. Combin. Theory \textbf{1} (1966), 209--214.

\bibitem{KorteLovaszSchrader91}
B.~Korte, L.~Lov\'asz, and R.~Schrader,
\emph{Greedoids}, Algorithms and Combinatorics, vol.~4,
Springer-Verlag, Berlin--Heidelberg, 1991.

\bibitem{Kozlov08}
D.~Kozlov,
\emph{Combinatorial Algebraic Topology},
Algorithms and Computation in Mathematics, vol.~21, Springer, 2008.

\bibitem{Roth96}
D.~Roth,
\emph{On the hardness of approximate reasoning},
Artificial Intelligence \textbf{82} (1996), 273--302.

\bibitem{SangEtAl04}
T.~Sang, F.~Bacchus, P.~Beame, H.~A. Kautz, and T.~Pitassi,
\emph{Combining component caching and clause learning for effective model counting},
in \emph{Online Proceedings of the Seventh International Conference on Theory and Applications of Satisfiability Testing (SAT 2004)},
2004, pp.~20--28.

\bibitem{XueChoiDarwiche12}
Y.~Xue, A.~Choi, and A.~Darwiche,
\emph{Basing decisions on sentences in decision diagrams},
in \emph{Proceedings of the Twenty-Sixth AAAI Conference on Artificial Intelligence},
AAAI Press, 2012, pp.~842--849.
\end{thebibliography}

\end{document}
